\section{Related Work}
\label{sec:related_work}

\subsection{Teleoperation and Demonstration Collection}
Low-cost isomorphic interfaces, master--slave systems, wearable hand capture,
and visual retargeting have reduced the hardware barrier to collecting
real-robot demonstrations and expanded the range of accessible embodiments
\citep{wuGELLOGeneralLowCost2024,zhaoLearningFineGrainedBimanual2023,fuMobileALOHALearning2024,chiUniversalManipulationInterface2024,wangDexCapScalablePortable2024,qinAnyTeleopGeneralVisionBased2024,xuDexUMIUsingHuman2025}.
DROID, Open X-Embodiment, EgoDex, and MimicGen further demonstrate paths toward
cross-scene data collection, egocentric data, and demonstration-driven
augmentation
\citep{khazatskyDROIDLargeScaleInTheWild2024,collaborationOpenXEmbodimentRobotic2025,hoqueEGODEXLEARNINGDEXTEROUSMANIPULATION2026,mandlekarMimicGenDataGeneration2023}.
These systems primarily treat operator commands as demonstrations to record.
Our method retains their teleoperation interface but treats the command stream
as a control signal whose filtering strength can change continuously with the
observed task state.

\subsection{Data Quality and Visual Auditing}
Demonstration-quality research uses information content, coverage, trajectory
properties, or task and model factors to decide which offline samples are worth
retaining
\citep{hejnaRobotDataCuration2025,chiDiffusionPolicyVisuomotor2024,zhaoLearningFineGrainedBimanual2023}.
This establishes the importance of downstream validation, but acts after the
operator time has already been spent. In parallel, visual and vision-language
models estimate success, failure, phase, progress, near misses, or process
reward
\citep{duanAHAVisionLanguageModelDetecting2024,sermanetRoboVQAMultimodalLongHorizon2024,leeRoboRewardGeneralPurposeVisionLanguage2026,wuLargeRewardModels2026}.
We use reviewed offline audit proposals to construct training views and data
admission, allowing audit results to affect subsequent collection rather than
only rank completed trajectories.

\subsection{Shared Control, Residual Policies, and Corrections}
Residual reinforcement learning, shared autonomy, deployment-time human
corrections, and interactive policy repair show that a base controller can be
retained while constrained corrections improve recovery
\citep{johanninkResidualReinforcementLearning2019,javdaniSharedAutonomyHindsight2015,liuRobotLearningJob2023,luoPreciseDexterousRobotic,welteFlowCorrectEfficientInteractive2026a}.
These works primarily evaluate task success, policy adaptation, or deployment
recovery. Our residual is instead a collection-time assistance signal: visual
and action history determine a candidate correction, a separate continuous
gain regulates its authority, and a safety projection constrains the composed
command. Nominal successes suppress unnecessary offsets, while correction
intervals supervise behavior around difficult states.

\subsection{Interactive Imitation and Iterative Collection}
DAgger, ThriftyDAgger, and deployment-time intervention learning address
distribution shift through expert queries, takeovers, or continued data
collection
\citep{rossReductionImitationLearning2011,hoqueThriftyDAggerBudgetAwareNovelty2021,liuRobotLearningJob2023,luoPreciseDexterousRobotic}.
They establish that collection, learning, and deployment can form a loop, but
usually optimize policy success or query cost. We instead measure the loop at
the data layer: under fixed operator time, reset count, configuration quotas,
and safety limits, how many newly collected episodes enter the action-training
set? Downstream ACT or Diffusion Policy performance is an independent utility
test, not an admission signal.

Together, prior work supplies teleoperation interfaces, offline curation,
visual auditing, shared control, and interactive learning. We connect these
capabilities through a collection-time continuous filter: audited history
supervises action and gain, real execution generates new observations, and a
fixed-budget protocol determines whether the next round produces more useful
data.
